\documentclass[11pt]{article}
\usepackage[a4paper,margin=25mm]{geometry}
\usepackage{amsmath,amssymb,amsthm,booktabs,array}
\usepackage[T1]{fontenc}
\usepackage[colorlinks=true,urlcolor=blue,citecolor=blue,linkcolor=blue]{hyperref}
\newtheorem{theorem}{Theorem}
\newtheorem{lemma}[theorem]{Lemma}
\newtheorem{proposition}[theorem]{Proposition}
\newtheorem{corollary}[theorem]{Corollary}
\newcommand{\DIO}{\operatorname{DIO}}
\newcommand{\Ad}{\operatorname{Ad}}
\newcommand{\Tr}{\operatorname{Tr}}
\newcommand{\id}{\operatorname{id}}
\title{Fixed points of superchannels for dephasing-covariant channels}
\author{Haobo Ma\\\small ChronoAI Pte Ltd, Singapore
\and Wenlin Zhang\\\small National University of Singapore\\\small ChronoAI Pte Ltd, Singapore}
\date{12 November 2026}
\begin{document}
\maketitle
\begin{abstract}
Extending the two-sided destruction construction of Saxena, Chitambar and
Gour (2019), we answer Yoeli and Gour's fixed-superchannel question in
\emph{Instability as a Quantum Resource},
arXiv:2602.18323v1, p.~5, Open directions, negatively:
in every dimension $d\geq3$, the dephasing-covariant channels are not exactly
the fixed channels of any deterministic superchannel.
Indeed, any superchannel fixing all such channels must be the identity on
the full space of linear maps. The proof uses the completely depolarizing
channel and a connected spanning family of monomial-unitary Choi vectors.
For qubits an explicit idempotent superchannel has the requested fixed set,
and the same construction works for pinching into two orthogonal sectors.
The obstruction also holds for uniformity-preserving channels in every
dimension at least two, and both obstructions persist with an unrestricted
quantum side system. We separate these results from the known two-sided
destruction construction, and exhibit the positivity obstruction to a
natural linear projection onto the covariant maps.
\end{abstract}

\section{The question and the setting}
Yoeli and Gour, in \emph{Instability as a Quantum Resource},
arXiv:2602.18323v1, p.~5, Open directions, identify as a ``key open question'' whether
``destruction-covariant channels themselves arise as fixed points of an
idempotent superchannel'' \cite{YG1}.
We answer this question negatively for complete dephasing in dimension
three and above. Complete dephasing is an admissible destruction channel:
it is idempotent, completely positive and trace preserving (CPTP), and fixes
the full-rank state $I/d$. The quoted question is absent from the revised
Open directions paragraph \cite[v2, p.~5]{YG2}; our posing locator is v1.

Let $M_d$ denote the complex $d\times d$ matrices, with standard basis
$|0\rangle,\ldots,|d-1\rangle$ and matrix units $E_{jk}=|j\rangle\langle k|$.
Write
\[
 \Delta_d(X)=\sum_{j=0}^{d-1}E_{jj}XE_{jj},\qquad
 \mathcal D_d(X)=\Tr(X)I/d,\qquad \Ad_U(X)=UXU^*.
\]
A channel $\mathcal N:M_d\to M_d$ is dephasing covariant if
$\mathcal N\Delta_d=\Delta_d\mathcal N$; denote this class by $\DIO_d$.
This is the commuting condition for the dephasing resource-destroying
map \cite{LHL}. It is distinct from covariance of a superchannel under
the two-sided channel dephasing construction of Saxena, Chitambar and
Gour (2019 preprint; 2020 publication) \cite{SCG}.
The latter framework takes classical channels as its fixed objects;
it does not identify all dephasing-covariant channels with fixed objects.
Ahnefeld, Theurer and Plenio \cite[Definition~9, Theorem~10]{ATP}
characterize compatible higher-order operations for DIO. Compatibility
under composition is different from the existence of a retraction fixing
every DIO channel, which is the question considered here.

Let $\mathsf L_d$ be the complex vector space of all complex-linear maps
$M_d\to M_d$. Use the input-first Choi convention
\[
 J_{\mathcal N}=\sum_{j,k}E_{jk}\otimes\mathcal N(E_{jk}),\qquad
 |U\rangle\!\rangle=\sum_j|j\rangle\otimes U|j\rangle.
\]
Then $J_{\Ad_U}=|U\rangle\!\rangle\langle\!\langle U|$ and
$J_{\mathcal D_d}=I_{d^2}/d$. A deterministic superchannel is a physical
pre- and post-processing transformation of channels, allowing a memory
system. Its complex-linear extension $\Theta:\mathsf L_d\to\mathsf L_d$
induces a completely positive map $T:M_{d^2}\to M_{d^2}$ through
$T(J_{\mathcal N})=J_{\Theta(\mathcal N)}$
\cite[pp.~2--3, Eqs.~(6)--(7) and Theorem~1]{CDP}.

\begin{theorem}\label{thm:rigidity}
Let $d\geq3$. If a complex-linear transformation
$\Theta:\mathsf L_d\to\mathsf L_d$ has completely positive Choi action
and fixes every channel in $\DIO_d$, then $\Theta=\id$ on $\mathsf L_d$.
Consequently, no deterministic superchannel has fixed CPTP set exactly
$\DIO_d$, even without requiring idempotence.
\end{theorem}

\section{Proof by fixed rank-one projections}
The projection-commutation step is a special case of
Arias, Gheondea and Gudder \cite[Corollary~3.4]{AGG}; we include its
elementary proof and the connected-family consequence.
\begin{lemma}\label{lem:rigidity}
Let $T:M_n\to M_n$ be completely positive and unital. Suppose that $T$
fixes the rank-one projections onto unit vectors $v_\alpha$ spanning
$\mathbb C^n$. Join two indices when
$\langle v_\alpha,v_\beta\rangle\ne0$. If this graph is connected, then
$T=\id$.
\end{lemma}
\begin{proof}
Choose a Kraus representation $T(X)=\sum_r A_rXA_r^*$.
For a fixed projection $P$, both $T(P)=P$ and $T(I-P)=I-P$ hold. Hence
\[
 0=(I-P)T(P)(I-P)
   =\sum_r ((I-P)A_rP)((I-P)A_rP)^*.
\]
Each summand is positive, so $(I-P)A_rP=0$ for every $r$.
Applying the same argument to $I-P$ gives $PA_r(I-P)=0$.
Thus $A_rP=PA_r$. Write $A_rv_\alpha=\lambda_{r,\alpha}v_\alpha$.
Commutation with $P_\alpha=|v_\alpha\rangle\langle v_\alpha|$ also gives
$\langle v_\alpha|A_r=\lambda_{r,\alpha}\langle v_\alpha|$.
Consequently
\[
 (\lambda_{r,\alpha}-\lambda_{r,\beta})
       \langle v_\alpha,v_\beta\rangle=0.
\]
Connectivity makes each eigenvalue independent of $\alpha$; spanning then
gives $A_r=\lambda_r I$. Finally, unitality implies
$\sum_r|\lambda_r|^2=1$, so $T(X)=X$.
\end{proof}

\begin{proof}[Proof of Theorem~\ref{thm:rigidity}]
The channel $\mathcal D_d$ commutes with $\Delta_d$, since both
$\Delta_d\mathcal D_d$ and $\mathcal D_d\Delta_d$ equal $\mathcal D_d$.
Its fixed Choi matrix $I_{d^2}/d$ implies $T(I_{d^2})=I_{d^2}$.

For a permutation $p$ of $\{0,\ldots,d-1\}$, let $P_p|j\rangle=|p(j)\rangle$.
Put $D_\varnothing=I$, and for each index $j$ let $D_j$ be diagonal with
$i$ at position $j$ and $1$ elsewhere, where $i^2=-1$.
All monomial unitaries $U_{p,q}=P_pD_q$, for
$q\in\{\varnothing,0,\ldots,d-1\}$, satisfy
$\Ad_{U_{p,q}}\Delta_d=\Delta_d\Ad_{U_{p,q}}$:
diagonal phases cancel on diagonal entries, and permutations relabel those
entries. Therefore $T$ fixes the rank-one projections onto
\[
 v_{p,q}=d^{-1/2}|U_{p,q}\rangle\!\rangle.
\]
These vectors span $\mathbb C^d\otimes\mathbb C^d$, because
\[
 v_{p,j}-v_{p,\varnothing}
   =\frac{i-1}{\sqrt d}|j\rangle\otimes|p(j)\rangle,
\]
and for each pair $j,k$ a permutation with $p(j)=k$ exists.

Each $v_{p,j}$ is joined to $v_{p,\varnothing}$ since
\[
 \langle v_{p,\varnothing},v_{p,j}\rangle=(d-1+i)/d\ne0.
\]
If $p'=p\circ\tau$ for a transposition $\tau$, then
\[
 \langle v_{p,\varnothing},v_{p',\varnothing}\rangle
   =\Tr(P_p^*P_{p'})/d=(d-2)/d\ne0.
\]
Every permutation is a product of transpositions, so the whole graph is
connected. Lemma~\ref{lem:rigidity} gives $T=\id$, hence $\Theta=\id$
by the Choi isomorphism.

To see why this forbids the specified fixed set, take
\[
 R=\begin{pmatrix}3/5&-4/5\\4/5&3/5\end{pmatrix}\oplus I_{d-2}.
\]
This is unitary, and $(RE_{00}R^*)_{01}=12/25$, whereas
$(\Delta_d(RE_{00}R^*))_{01}=0$. Thus $\Ad_R$ is not in $\DIO_d$,
but the identity superchannel fixes it.
\end{proof}

\section{Positive cases and the nonuniformity obstruction}
\begin{corollary}\label{prop:qubit}
Let $P$ be an orthogonal projection on $\mathbb C^d$,
$Q=I-P$, and $\Lambda(X)=PXP+QXQ$.
There is an idempotent deterministic superchannel whose fixed linear maps
are exactly those commuting with $\Lambda$.
In particular, the fixed-set question for complete dephasing has a positive
answer in dimensions $1$ and $2$.
\end{corollary}
\begin{proof}
Let $Z=P-Q$, so $Z^*=Z$, $Z^2=I$. Put $\mathcal C=\Ad_Z$ and define
\[
 \Gamma(\mathcal N)=\mathcal C\mathcal N\mathcal C,
 \qquad \Theta=(\id+\Gamma)/2.
\]
Both terms are physical pre-/post-processings, and their equal random
mixture is a deterministic superchannel. Since $\mathcal C^2=\id$,
we have $\Gamma^2=\id$ and $\Theta^2=\Theta$. Moreover,
\[
 \Theta(\mathcal N)=\mathcal N
 \ \Longleftrightarrow\ \mathcal C\mathcal N\mathcal C=\mathcal N
 \ \Longleftrightarrow\ \mathcal N\mathcal C=\mathcal C\mathcal N
 \ \Longleftrightarrow\ \mathcal N\Lambda=\Lambda\mathcal N,
\]
where the last step uses $\Lambda=(\id+\mathcal C)/2$.
For $d=2$ and $P=E_{00}$, $\Lambda=\Delta_2$.
For $d=1$, all linear maps commute with $\Delta_1=\id$, and the identity
superchannel supplies the corresponding fixed set.
\end{proof}

\begin{corollary}\label{cor:uniform}
Let $a,b\geq1$, and consider channels on $M_a\otimes M_b$.
For either destruction $\mathcal D_a\otimes\id_b$ with $a\geq2$, or
$\Delta_a\otimes\id_b$ with $a\geq3$, any complex-linear transformation
with completely positive Choi action fixing every destruction-covariant
channel is the identity. Neither covariant class is an exact physical
fixed set. For $\Delta_a\otimes\id_b$, such a fixed-set projection exists
exactly when $a\leq2$.
In particular, the answer is negative for nonuniformity and conditional
nonuniformity in Table~I of Yoeli and Gour.
\end{corollary}
\begin{proof}
First take $b=1$, write $d=a$, and consider destruction $\mathcal D_d$.
For CPTP $\mathcal N$, commuting with $\mathcal D_d$ is equivalent to
$\mathcal N(I/d)=I/d$. This class contains $\mathcal D_d$ and every
unitary channel. For $d\geq3$ the proof of
Theorem~\ref{thm:rigidity} therefore applies unchanged.
For $d=2$, the six monomial vectors used there split into two connected
families, for the identity and transposition permutations. Their span is
still all of $\mathbb C^4$. The rational rotation $R$ displayed above
connects them: with $X=\bigl(\begin{smallmatrix}0&1\\1&0\end{smallmatrix}\bigr)$
and $D_0=\operatorname{diag}(i,1)$,
\[
 \tfrac12\Tr(R)=3/5\ne0,\qquad
 \tfrac12\Tr(R^*XD_0)=2(i-1)/5\ne0.
\]
The seven-vector family is connected, so Lemma~\ref{lem:rigidity} applies.
Finally, the channel $\mathcal N_0(X)=\Tr(X)E_{00}$ is not
uniformity preserving. The identity superchannel fixes it, excluding the
requested exact fixed set.

For the tensor extensions, in each dimension $r$ choose the connected
spanning unitary Choi family just constructed: the monomial family for
$r\geq3$, the seven-vector family for $r=2$, and the single identity
vector for $r=1$. For dephasing on the first factor with $a\geq3$, use
only its monomial family there. Every product unitary $U_A\otimes V_B$
so obtained is covariant for the specified destruction. Its normalized
Choi vector is the tensor product of the two local vectors, up to a fixed
reordering of tensor factors. These vectors span the full Choi space, and
their overlap graph is connected: an edge in either local graph gives an
edge with the other coordinate fixed. Both destructions are unital, so
their covariant classes contain the global completely depolarizing channel.
The Choi action is therefore unital, and Lemma~\ref{lem:rigidity} applies.
For $\mathcal D_a\otimes\id_b$, the channel
$\mathcal N_0\otimes\id_b$ is not covariant; for
$\Delta_a\otimes\id_b$, use $\Ad_R\otimes\id_b$ instead.
Finally, $a=2$ is two-sector pinching with $P=E_{00}\otimes I_b$,
covered by Corollary~\ref{prop:qubit}; $a=1$ gives the identity destruction.
\end{proof}

\section{Two-sided destruction and linear projections}
Let $\Delta_A:M_{d_A}\to M_{d_A}$ and
$\Delta_B:M_{d_B}\to M_{d_B}$ be idempotent destruction channels,
and let $\mathcal N:M_{d_A}\to M_{d_B}$.
The known
construction $\Pi(\mathcal N)=\Delta_B\mathcal N\Delta_A$
is physical and idempotent. Its fixed maps satisfy exactly
$\mathcal N=\Delta_B\mathcal N=\mathcal N\Delta_A$:
apply either idempotent to $\mathcal N=\Pi(\mathcal N)$; the converse is
immediate. These are stronger conditions than
$\Delta_B\mathcal N=\mathcal N\Delta_A$.
For complete dephasing, this is the two-sided construction of
Saxena, Chitambar and Gour (2019 preprint; 2020 publication)
\cite[Eq.~(4)]{SCG}, reproduced in
Eq.~(29) of Yoeli--Gour v1 and Eq.~(31) of v2.
In the equal-dimensional case it is
$\Pi_d(\mathcal N)=\Delta_d\mathcal N\Delta_d$
\cite[p.~4]{YG1}; it becomes Eq.~(31) in v2 \cite[p.~5]{YG2}.
It is physical and idempotent. Its fixed channels are precisely
\[
 \mathcal N(X)=\sum_{j,k}s_{k|j}X_{jj}E_{kk},
 \qquad s_{k|j}\geq0,\quad\sum_k s_{k|j}=1.
\]
Indeed, $\mathcal N=\Pi_d(\mathcal N)$ removes every off-diagonal input
matrix unit and forces each $\mathcal N(E_{jj})$ to be diagonal;
positivity and trace preservation give the stated stochastic columns.
Conversely, every displayed channel is fixed by $\Pi_d$.
Such channels commute with $\Delta_d$, but for $d\geq2$ this inclusion is
strict: the identity channel commutes with $\Delta_d$ and is not fixed
by $\Pi_d$. Thus the obstruction concerns the larger covariant class,
not the validity of Eq.~(29).

The other explicit example following Eq.~(29) also survives unchanged.
More generally, put $\Delta_A=\id$ and
$\Delta_B(X)=\Tr(X)\gamma$ for a full-rank density matrix $\gamma$.
For CPTP $\mathcal N$, $\Delta_B\mathcal N$ is the constant channel
$X\mapsto\Tr(X)\gamma$. Hence covariance is precisely equality with that
channel, and $\mathcal N\mapsto\Delta_B\mathcal N$ is the required
idempotent superchannel. The source's uniform channel is $\gamma=I/d_B$.

\begin{corollary}\label{cor:linear}
For arbitrary idempotent linear maps $\Delta_A,\Delta_B$, the formula
\[
 \mathcal Q(\mathcal N)=\Delta_B\mathcal N\Delta_A
       +(\id-\Delta_B)\mathcal N(\id-\Delta_A)
\]
is an idempotent linear projection with fixed space exactly
$\{\mathcal N:\Delta_B\mathcal N=\mathcal N\Delta_A\}$.
For complete dephasing on $M_3$, it sends some CPTP maps to maps that
are not completely positive.
\end{corollary}
\begin{proof}
The idempotents split the input and output spaces into their ranges and
kernels. Covariance says that the two off-diagonal blocks of
$\mathcal N$ vanish, and $\mathcal Q$ retains exactly the two diagonal
blocks. This proves idempotence and the fixed-space characterization.
For the positivity claim take the real unitary
\[
 U=\frac19\begin{pmatrix}7&-4&-4\\-4&1&-8\\-4&-8&1\end{pmatrix}
   =I-\frac29\begin{pmatrix}1\\2\\2\end{pmatrix}
                  \begin{pmatrix}1&2&2\end{pmatrix}.
\]
In the Choi matrix of $\mathcal Q(\Ad_U)$, the principal submatrix indexed
by $(0,0),(1,1),(0,2)$ is
\[
 \frac1{81}\begin{pmatrix}49&7&0\\7&1&-4\\0&-4&16\end{pmatrix}.
\]
Its determinant is $-784/531441<0$, so the Choi matrix is not positive.
This is an obstruction to this particular linear projection; the
nonexistence of \emph{any} physical projection follows from
Theorem~\ref{thm:rigidity}.
\end{proof}

The same issue already appears for affine sets of states.
For completely positive resource-destroying maps, Gour
\cite[Appendix~F, Theorem~2]{Gour2017} already requires Choi positivity
in addition to affinity. The following example concerns the weaker
positivity requirement in the state-set definition below.
In the terminology of Gour \cite[Definition~9.7]{Gour}, a set is affine
when any affine combination of its elements that is a density matrix
remains in the set.
\begin{corollary}\label{cor:affine}
The affine set $F=\{\rho\in M_3:\rho\geq0,\ \Tr\rho=1,\ \rho_{02}=0\}$
admits no self-adjoint linear resource-destroying map, that is, no
Hilbert--Schmidt self-adjoint linear map taking every state into $F$
and fixing every state in $F$.
\end{corollary}
\begin{proof}
The real span of $F$ is the Hermitian subspace
$S=\{H=H^*:H_{02}=0\}$. Indeed, $I/3$ belongs to $F$, and every traceless
direction in $S$ has sufficiently small positive and negative perturbations
of $I/3$ inside $F$. Such a resource-destroying map would be an idempotent
with range $S$ on the Hermitian matrices. Self-adjointness makes it the
orthogonal projection onto $S$, which deletes entries $(0,2)$ and $(2,0)$.
Applied to the state $\frac13(1,1,1)^{\mathsf T}(1,1,1)$, it yields
\[
 \frac13\begin{pmatrix}1&1&0\\1&1&1\\0&1&1\end{pmatrix},
 \qquad\det=-1/27<0,
\]
which is not a state, a contradiction.
\end{proof}
Thus the assertion that ``there exists a unique self-adjoint RDM'' for
every affine set in Theorem~9.2 of \emph{arXiv:2402.05474v2},
printed p.~442, fails as stated.
This observation concerns that version, not an identification of its
theorem numbering or wording with a later book edition.

\paragraph{Covariance under a twirl versus its group.}
The distinction in v1, p.~4, Open directions, is already visible for
qubits. For the diagonal-phase group the twirl is $\Delta_2$.
The swap channel $\Ad_X$ commutes with $\Delta_2$, but does not commute
with $\Ad_{\operatorname{diag}(1,i)}$: on $E_{01}$ the two orders give
$-iE_{10}$ and $iE_{10}$. This reproves the source's strict distinction;
it does not settle its broader question about restricted free operations.

\section{Source coverage and verification scope}
\begin{center}\small
\begin{tabular}{@{}>{\raggedright\arraybackslash}p{0.43\linewidth}>{\raggedright\arraybackslash}p{0.51\linewidth}@{}}
\toprule
Source item & Mathematical coverage\\
\midrule
v1, p.~5, Open directions, fixed-superchannel question & Answered: no in general;
complete dephasing for all $d\geq3$,
Theorem~\ref{thm:rigidity}.\\
Complete-dephasing boundary cases & Not posed separately; proved for $d=1,2$;
Corollary~\ref{prop:qubit}.\\
Two-sector pinching; v1, Table~I, nonuniformity and conditional nonuniformity & Not posed separately;
positive construction and obstruction proved,
Corollaries~\ref{prop:qubit} and \ref{cor:uniform}.\\
Dephasing with an unrestricted side system & Not posed separately; same
classification by the dephased dimension, Corollary~\ref{cor:uniform}.\\
v1, Eq.~(29); v2, Eq.~(31), classical channels & Reproved; prior proof
Saxena--Chitambar--Gour (2019/2020), Eq.~(4).\\
v1, p.~4, after Eq.~(29), uniform channel & Reproved; prior example
Yoeli--Gour (2026); extended to full-rank replacers.\\
v1, p.~4, group twirl versus group covariance & Reproved; prior distinction
Yoeli--Gour (2026); explicit qubit witness above.\\
Linear projection onto the covariance subspace & Not posed; obstruction proved,
Corollary~\ref{cor:linear}.\\
v1 reference [26], arXiv:2402.05474v2, Theorem~9.2 & Refuted (explicit
counterexample), Corollary~\ref{cor:affine}.\\
v1, Open directions: nonphysical destruction maps and restricted free channels
& Not covered.\\
v1, Open directions: higher dynamical/GPT reversibility
& Not covered.\\
\bottomrule
\end{tabular}
\end{center}

In the Lean module \texttt{YoeliGourDestructionSuperchannelRefutation},
the pinned declaration
\href{https://github.com/the-omega-institute/trureturing/blob/a252cc46f0e145418a978217b8abd41a73fd820d/D5/S3/QuantumChannels/YoeliGourDestructionSuperchannelRefutation.lean}{\texttt{result}}
proves the negation of the universal admissible-system fixed-set claim,
using qutrit complete dephasing. Its helper
\texttt{fixes\_dio\_identity} establishes the qutrit rigidity statement.
All other results above have written proofs, including the dimension-uniform
extension, the boundary cases and the sibling corollaries. The accompanying
exact-arithmetic program checks finite qutrit and qubit families, selected
two-sector pinchings, both negative determinants, the replacer identity and
the group-covariance witness. These finite checks do not supply an
all-dimension proof.

\begin{thebibliography}{9}
\bibitem{YG1} G.~Yoeli and G.~Gour,
\emph{Instability as a Quantum Resource},
\href{https://arxiv.org/abs/2602.18323v1}{arXiv:2602.18323v1} (2026),
especially pp.~1, 4--5, Eq.~(1), Eq.~(29), and Open directions.
\bibitem{YG2} G.~Yoeli and G.~Gour,
\emph{Instability as a Quantum Resource},
\href{https://arxiv.org/abs/2602.18323v2}{arXiv:2602.18323v2} (2026),
especially p.~5, Eq.~(31) and Open directions.
\bibitem{LHL} Z.-W.~Liu, X.~Hu, and S.~Lloyd,
\emph{Resource Destroying Maps}, Phys. Rev. Lett. \textbf{118}, 060502 (2017),
\href{https://doi.org/10.1103/PhysRevLett.118.060502}{doi:10.1103/PhysRevLett.118.060502}.
\bibitem{CDP} G.~Chiribella, G.~M.~D'Ariano, and P.~Perinotti,
\emph{Transforming quantum operations: quantum supermaps},
EPL \textbf{83}, 30004 (2008),
\href{https://doi.org/10.1209/0295-5075/83/30004}{doi:10.1209/0295-5075/83/30004}.
\bibitem{SCG} G.~Saxena, E.~Chitambar, and G.~Gour,
\emph{Dynamical Resource Theory of Quantum Coherence},
Phys. Rev. Research \textbf{2}, 023298 (2020),
\href{https://doi.org/10.1103/PhysRevResearch.2.023298}{doi:10.1103/PhysRevResearch.2.023298}.
\bibitem{ATP} F.~Ahnefeld, T.~Theurer, and M.~B.~Plenio,
\emph{The Structure of Higher-Order Quantum Resources with an Application
to Coherence},
\href{https://arxiv.org/abs/2609.40186v1}{arXiv:2609.40186v1} (2026),
Definition~9, Theorem~10 and Appendix~C, Eq.~(C252c).
\bibitem{Gour} G.~Gour, \emph{Resources of the Quantum World},
\href{https://arxiv.org/abs/2402.05474v2}{arXiv:2402.05474v2} (2024),
Section~9.3.1, Definitions~9.9--9.10 and Theorem~9.2,
printed pp.~440--443.
\bibitem{AGG} A.~Arias, A.~Gheondea, and S.~Gudder,
\emph{Fixed points of quantum operations},
J. Math. Phys. \textbf{43}, 5872--5881 (2002),
\href{https://doi.org/10.1063/1.1519669}{doi:10.1063/1.1519669},
Corollary~3.4.
\bibitem{Gour2017} G.~Gour,
\emph{Quantum resource theories in the single-shot regime},
Phys. Rev. A \textbf{95}, 062314 (2017),
\href{https://arxiv.org/abs/1610.04247v3}{arXiv:1610.04247v3},
Appendix~F, Theorem~2.
\end{thebibliography}
\bigskip\noindent\small
\copyright\ 2026 Haobo Ma and Wenlin Zhang.
Licensed under Creative Commons Attribution 4.0 International
(\href{https://creativecommons.org/licenses/by/4.0/}{CC BY 4.0}).
\end{document}
